%% ===================================================================================
%% Apêndice E -- Registros de Decisão Arquitetural (ADR)
%% Formato: Nygard (2011). Uma decisão aceita só muda por uma NOVA ADR que a substitui
%% (status "Substituída por ADR-0xx"); nunca se apaga uma ADR. Questões ainda abertas
%% ficam no Apêndice F. Dono: Arquiteto/Redator.
%% ===================================================================================
\chapter{Registros de Decisão Arquitetural}\label{ap:adr}
\begingroup\small

Este apêndice reúne as decisões arquiteturais aceitas, no formato de \en{Architecture Decision Records} \cite{nygard2011adr}: contexto, decisão, consequências e alternativas. Uma decisão só é alterada por uma nova ADR que a substitua; as questões ainda não decididas estão no Apêndice~\ref{ap:decisoes}.

\adrentry{ADR-001}{Monólito modular com Next.js (App Router)}{Aceita.}
{O Vitrine é um sistema de um único \en{tenant}, com tráfego moderado, implantado pelo próprio docente; a equipe é pequena e há agentes de desenvolvimento em paralelo. O requisito de implantação simples (\req{RNF01}) e o baixo custo operacional são prioritários.}
{Adotar uma única aplicação Next.js (App Router), organizada internamente em módulos por domínio --- Dados, Core e UI --- com dependências unidirecionais (Figuras~\ref{fig:arq-componentes} e~\ref{fig:arq-camadas}).}
{Positivas: implantação e depuração simples, transações locais, uma única unidade de versionamento. Negativas: escalabilidade horizontal limitada (cache ISR por instância, SQLite em arquivo); a interface e a API são implantadas juntas.}
{\en{Front-end} separado com API independente; microsserviços; CMS \en{headless} de terceiros. Descartadas pelo custo de operação e de coordenação desproporcional ao escopo \cite{newman2019monolith,fowler2015monolithfirst}.}

\adrentry{ADR-002}{SQLite como banco de dados embutido}{Aceita.}
{O objetivo OE4 pede uma arquitetura de dados portátil, sem servidor de banco de dados, com baixo custo para o usuário final. O perfil de uso é de leitura predominante, com poucas escritas realizadas por um administrador.}
{Usar o SQLite, com o banco em um único arquivo em volume persistente e, preferencialmente, em modo WAL.}
{Positivas: sem servidor a administrar, cópia de segurança simples, custo zero. Negativas: escrita serializada; necessidade de disco persistente (incompatível com sistemas de arquivos efêmeros, risco R01, tratado pela \adr{ADR-011}); ausência de acesso remoto nativo; valores de enumeração não impostos pelo banco \cite{sqliteabout,sqlitewal}.}
{PostgreSQL ou MySQL (exigem servidor); libSQL/Turso (banco remoto compatível com SQLite, via possível para plataformas \emph{serverless}, descartada nesta versão pela \adr{ADR-011}).}

\adrentry{ADR-003}{Prisma ORM 7 como camada de acesso a dados}{Aceita.}
{É necessário um modelo de dados declarativo, versionável e com migrações, e um cliente consistente para o SQLite. O projeto usa o Prisma 7, cujas práticas diferem das versões anteriores (configuração em arquivo próprio, novo gerador, \en{driver adapter}).}
{Adotar o Prisma ORM 7. Somente \arq{src/lib/db.js} instancia o cliente; o restante do código usa as funções exportadas por ele.}
{Positivas: esquema único como fonte do modelo, migrações e cliente tipado. Negativas: a adequação à versão 7 gerou trabalho inicial, resolvido na \adr{ADR-014}; o Prisma acrescenta uma camada de abstração entre a aplicação e o SQL.}
{Drizzle ORM; SQL direto com um \en{driver} de SQLite (mais controle, menos produtividade); uso apenas do \en{query builder} interno do better-auth.}

\adrentry{ADR-004}{better-auth para autenticação, em substituição ao Auth.js}{Aceita.}
{O relatório de IC previu o Auth.js \cite{reis2026ic}. Na consolidação do monólito, o repositório adotou o better-auth: a dependência está declarada, o esquema de dados contém as tabelas \texttt{user}, \texttt{session}, \texttt{account} e \texttt{verification} no formato do better-auth, e o pacote instalado oferece adaptador para Prisma e integração com Next.js.}
{Adotar o better-auth 1.7.x com adaptador do Prisma, autenticação por e-mail e senha, sessões persistidas em banco e cadastro público desabilitado após o \en{bootstrap}.}
{Positivas: esquema alinhado ao que já existe; integração com Next.js; \en{plugins} para evolução (papéis, autenticação de dois fatores). Negativas: biblioteca mais recente, exigindo acompanhar mudanças de API; todas as menções ao Auth.js do relatório de IC passam a ser históricas.}
{Auth.js (prevista na IC); implementação própria, descartada pelo risco de segurança. \textbf{Motivação original a confirmar com o autor}: esta ADR registra o estado do repositório e justificativas técnicas verificáveis; a motivação da troca deve ser completada pelo autor.}

\adrentry{ADR-005}{ISR com revalidação sob demanda}{Aceita.}
{As páginas públicas precisam de desempenho de página estática (\req{RNF02}) sem exigir novo \en{build} a cada alteração de conteúdo ou de tema (OE3).}
{Pré-renderizar as páginas públicas com \arq{generateStaticParams}; definir \texttt{revalidate} alto como rede de segurança; invalidar sob demanda com \arq{revalidatePath} nas mutações (regra RN10); não habilitar Cache Components.}
{Positivas: desempenho e atualização imediata na próxima visita. Negativas: suporta apenas o \en{runtime} Node.js; o cache é por instância (exige instância única); a invalidação do \en{layout} raiz regenera todas as páginas \cite{nextjsdocs}.}
{SSR a cada requisição; SSG puro com novo \en{build}; \texttt{revalidateTag} com granularidade por etiqueta (reavaliar se o volume de conteúdo crescer); Cache Components.}

\adrentry{ADR-006}{Markdown com \texttt{unified}, KaTeX e sem HTML bruto}{Aceita.}
{O objetivo OE2 exige escrita científica: fórmulas \LaTeX{}, código com realce e tabelas. A segurança exige que o conteúdo do administrador nunca execute \en{scripts} no navegador do leitor.}
{Renderizar com \arq{react-markdown}, \arq{remark-math} e \arq{rehype-katex}, acrescentando GFM e realce de sintaxe; proibir HTML bruto.}
{Positivas: texto legível e portátil; segurança por padrão; matemática nativa. Negativas: sem HTML arbitrário; dependência de \en{plugins} do ecossistema \texttt{unified} \cite{unifiedjs,reactmarkdown,katex}.}
{MDX (permite código executável, risco de segurança); editor visual que gera HTML; renderização apenas no cliente.}

\adrentry{ADR-007}{Tematização por variáveis CSS e Tailwind CSS 4}{Aceita.}
{A personalização visual deve refletir-se sem novo \en{build} (OE3).}
{Mapear tokens do Tailwind para variáveis CSS (\texttt{@theme inline}) e definir os valores no atributo \texttt{style} do elemento raiz, a partir da configuração do banco; validar cores como \texttt{\#RRGGBB}.}
{Positivas: mudança de tema sem recompilar CSS; baixa complexidade. Negativas: exige disciplina para usar sempre os tokens; convivência com o modo escuro é uma questão em aberto (\decisao{D-17}) \cite{tailwinddocs,mdncssvars}.}
{CSS-in-JS em tempo de execução; gerar CSS no \en{build} (exigiria recompilar); temas fixos predefinidos.}

\adrentry{ADR-008}{JavaScript com JSDoc, sem TypeScript}{Aceita.}
{O projeto foi iniciado em JavaScript (presença de \arq{jsconfig.json}) e precisa minimizar a fricção entre agentes.}
{Manter JavaScript (módulos ES), documentando as funções públicas de \arq{src/lib} com JSDoc. Reavaliar TypeScript em versões futuras.}
{Positivas: menor fricção e configuração. Negativas: menos garantias estáticas; erros de tipo aparecem em testes ou em execução, o que aumenta a importância dos testes de unidade e de integração \cite{bierman2014typescript,andreasen2017survey}.}
{TypeScript desde já.}

\adrentry{ADR-009}{Documentation-Driven Development com \LaTeX{} e diagramas como código}{Aceita.}
{O desenvolvimento por agentes em paralelo exige um contrato único e versionado; o trabalho é um TCC e a documentação também é o produto acadêmico.}
{Manter a documentação oficial em \arq{doc/tex} (\LaTeX{}, base abntex2) e os diagramas em \arq{doc/diagrams/src} (PlantUML e DBML), com um único redator (Arquiteto), figuras com \en{placeholder} automático e mapa de pendências.}
{Positivas: documentação revisável como código e pronta para a entrega; coerência entre agentes. Negativas: custo de redação e de manutenção pelo Arquiteto; dependência de ferramentas de composição e de renderização \cite{parnas1986rational,writethedocs}.}
{Documentação em \en{wiki} ou Markdown solto; documentos de texto proprietários; documentar só ao final.}

\adrentry{ADR-010}{Fluxo Git por domínios com mural na \texttt{main}}{Aceita.}
{Vários agentes trabalham simultaneamente e precisam de um canal de coordenação e de um ponto único de integração. O repositório não possui, atualmente, um remoto configurado.}
{Desenvolver apenas em ramificações por domínio; proibir código na \texttt{main}; usar \arq{MURAL_DE_AGENTES.md}, editado somente na \texttt{main}, como canal de pedidos de integração e de comunicação; o Arquiteto revisa, integra ou veta.}
{Positivas: rastreabilidade no histórico do Git; conflitos de comunicação evitados; integração controlada. Negativas: o Arquiteto é ponto único de revisão; trocar de ramificação exige árvore de trabalho limpa \cite{chacon2014progit,conventionalcommits}.}
{Pedidos de integração em plataforma Git hospedada (sem remoto hoje); trabalho direto na \texttt{main}.}

\adrentry{ADR-011}{Implantação por imagem Docker em VPS tradicional, com volume persistente}{Aceita (decisão do Engenheiro Chefe em 30/09/2026, pedido P-001; resolve \decisao{D-04}).}
{O SQLite exige disco persistente (\adr{ADR-002}), e o sistema de arquivos de plataformas \emph{serverless} é efêmero (risco R01); o relatório de IC citava a Vercel como exemplo de implantação \cite{reis2026ic}. O Implantador é, em geral, o próprio docente com apoio técnico, e a ISR exige \en{runtime} Node.js e instância única (\adr{ADR-005}).}
{Distribuir o Vitrine como \textbf{imagem Docker} executada em uma \textbf{VPS tradicional} (Linux com Docker Engine). O arquivo SQLite reside em um \textbf{volume persistente} montado no contêiner (por exemplo, em \arq{/data}); executa-se uma única instância; um \en{proxy} reverso fornece TLS; os segredos entram por variáveis de ambiente, fora da imagem; as cópias de segurança usam \texttt{VACUUM INTO}. Plataformas \emph{serverless} não são alvo desta versão.}
{Positivas: persistência garantida do banco; ambiente reprodutível e documentável por \texttt{Dockerfile} e \texttt{docker-compose.yml}; coerência com a ISR e com o modelo de instância única \cite{dockerdocs}. Negativas: exige uma VPS e noções básicas de Docker pelo Implantador (o requisito \req{RNF01} é preservado por poucos comandos de \en{compose}); o módulo nativo \arq{better-sqlite3} precisa ser instalado no \en{build} da imagem; o volume e as cópias de segurança são responsabilidade do operador; a CLI do Prisma (hoje em \texttt{devDependencies}) é necessária para aplicar migrações, o que deve ser resolvido no módulo M13.}
{Plataformas \emph{serverless} (Vercel) com banco remoto compatível com SQLite, como libSQL/Turso (descartada nesta versão: troca do \en{driver}, serviço externo e maior complexidade); PaaS com disco persistente; instalação direta com Node.js sem contêiner (continua possível, mas não é o alvo documentado).}

\adrentry{ADR-012}{Estado de publicação do conteúdo: DRAFT e PUBLISHED, com padrão PUBLISHED}{Aceita (decisão do Engenheiro Chefe em 30/09/2026; resolve \decisao{D-15}; integrada pelo merge M-001).}
{Os protótipos da IC trazem o comando ``Enviar nova versão'' e a possibilidade de preparar conteúdo antes de exibi-lo. A recomendação inicial do Arquiteto era a publicação imediata, com rascunho apenas local no navegador (\req{RF15}).}
{Adicionar o campo \texttt{Content.status}, do tipo enumerado \texttt{ContentStatus} (\texttt{DRAFT} ou \texttt{PUBLISHED}), com padrão \texttt{PUBLISHED}. Um conteúdo é visível ao público somente se ele e todos os seus ancestrais estiverem publicados (regra RN15). A API aceita e devolve o estado; alterá-lo é uma mudança estrutural para efeito de revalidação (RN10 e RN16). O histórico de versões continua fora do escopo.}
{Positivas: permite preparar e revisar conteúdo no servidor sem expô-lo. Negativas: toda leitura pública deve filtrar por estado (risco de vazamento de rascunhos, coberto pelo caso de teste CT31); a geração estática e o mapa do site devem excluir rascunhos; como o padrão é \texttt{PUBLISHED}, um conteúdo criado sem escolha explícita fica público, e a interface deve tornar o estado visível; o SQLite não impõe o enumerado no banco, que é validado pelo Prisma e pela aplicação.}
{Publicação imediata sem campo de estado (recomendação inicial); estado mais tabela de versões (descartada por exceder o escopo).}

\adrentry{ADR-013}{Links sociais do rodapé como campos individuais opcionais}{Aceita (decisão do Engenheiro Chefe em 30/09/2026; resolve \decisao{D-19}; integrada pelo merge M-001).}
{O rodapé dos protótipos exibe ícones de redes sociais, mas o esquema original não previa onde guardar os endereços. A recomendação inicial era um único campo de texto em Markdown.}
{Acrescentar a \texttt{Configuration} seis campos opcionais e independentes: \arq{discordUrl}, \arq{githubUrl}, \arq{youtubeUrl}, \arq{instagramUrl}, \arq{linkedinUrl} e \arq{twitchUrl}. Apenas URLs absolutas com protocolo \texttt{http} ou \texttt{https} são aceitas (RN17); campo vazio significa ícone omitido. O texto de direitos do rodapé é gerado a partir do nome do site e do ano corrente, sem campo próprio nesta versão.}
{Positivas: formulário simples e validável por campo; ícones padronizados; ausência de Markdown arbitrário no rodapé, o que reduz a superfície de ataque. Negativas: a lista de redes é fixa, e incluir outra rede exige migração e ajuste da interface.}
{Campo único em Markdown (recomendação inicial); campo JSON com a lista de links; tabela de links extensível (descartada por excesso de estrutura para a versão inicial).}

\adrentry{ADR-014}{Driver adapter better-sqlite3 e conformidade com o Prisma 7}{Aceita (proposta pelo agente Dados no pedido M-001 e aprovada pelo Arquiteto em 30/09/2026; resolve \decisao{D-09}, itens 1 a 5).}
{O Prisma 7 exige configuração da CLI em \arq{prisma.config.ts}, o gerador \texttt{prisma-client} e um \en{driver adapter} para o SQLite; a configuração inicial do repositório seguia práticas do Prisma 6 e declarava o pacote \arq{sqlite3}, que não é adaptador do Prisma. Com a implantação em VPS (\adr{ADR-011}), não há necessidade de banco remoto.}
{Adotar o adaptador \arq{@prisma/adapter-better-sqlite3} com o \arq{better-sqlite3} (série 12, em \texttt{dependencies}); renomear a configuração para \arq{prisma.config.ts}; usar o gerador \texttt{prisma-client} com saída em \arq{src/generated/prisma} (não versionada); retirar a URL do \texttt{datasource} do esquema; declarar o \arq{dotenv} em \texttt{devDependencies} para a CLI; mover \arq{@prisma/client} para \texttt{dependencies}; remover o \arq{sqlite3}.}
{Positivas: acesso direto ao arquivo SQLite, sem rede, coerente com a implantação em contêiner. Negativas: \arq{better-sqlite3} é módulo nativo (o \en{script} de instalação precisa ser permitido em \texttt{allowScripts}, e o \en{build} da imagem precisa de ferramentas de compilação ou binários pré-compilados); o cliente gerado é escrito em TypeScript dentro de um projeto JavaScript, e sua compilação pelo Next.js deve ser validada ao implementar o \arq{db.js}; \arq{prisma} e \arq{dotenv} permanecem em \texttt{devDependencies}, o que afeta a aplicação de migrações na imagem de produção (módulo M13).}
{\arq{@prisma/adapter-libsql} com \arq{@libsql/client} (banco remoto ou \en{edge}; mantido como via futura, não adotada pela \adr{ADR-011}); o pacote \arq{sqlite3} (não é adaptador do Prisma).}

\endgroup

